\color{unidarkgreen}{\subsection{Логика отключения}}
\color{black}

\begin{enumerate}


%%%%%%%%%%%%%%%%%%%%%%%%%%%%%%%%%%%%%%%%%%%%%%%%%% Логика запрета АПВ %%%%%%%%%%%%%%%%%%%%%%%%%%%%%%%%%%%%%%%%%%%%%%%%%%
\item Логика запрета АПВ

\begin{enumerate}

\item
Логика запрета АПВ предназначена для формирования запрета АПВ всех присоединений СШ, чтобы минимизировать вероятность повторного возникновения аварийной ситуации. 

\item
В устройстве предусмотрена реализация два типа АПВ: общий запрет АПВ и избирательный запрет АПВ при срабатывании ДЗШ.

\item
Сигналы общего запрета АПВ формируется в следующих условиях:
\begin{itemize}
\item при неуспешном АПВ;
\item при неполнофазном или полнофазном отказе выключателя питающего присоединения;
\item при срабатывании УРОВ;
\item при неуспешном опробовании;
\item при внешнем отключении;
\item при оперативном запрете АПВ.
\end{itemize}

\item
Неуспешное АПВ выявляется при повторном срабатывании защиты после первого отключения.
Для фиксации первого отключения используется сигнал \textit{«Фикс. откл. 1С»} (\textit{«Фикс. откл. 2С»}), который формируется в логике отключения.
Для того чтобы сигнал неуспешного АПВ формировался только при повторном срабатывании защиты, используется дополнительная выдержка времени, задаваемая уставкой \textbf{«Т АПВ 1пр. 1С»} (\textbf{«Т АПВ 1пр. 2С»}).
Данная выдержка времени должна охватывать время отключения всех выключателей, но при этом должна быть отстроена от времени включения первого присоединения в цикле АПВ.

\item
Неполнофазный или полнофазный отказ выключателя питающего присоединения выявляется по факту наличия напряжения после сформированного сигнала отключения.
Отключение СШ фиксируется при наличии сигнала \textit{«Фикс. откл. 1С»} (\textit{«Фикс. откл. 2С»}) и отсутствии сигнала \textit{«Отключение 1С»} (\textit{«Отключение 2С»}).
Выдержка времени 25 мс необходима для отстройки от разновременности отключения фаз выключателя.

\item
Формирование сигналов запрета АПВ при внешнем отключении описан {\color{red} в главе «Внешнее отключение».}

\item
Запрет АПВ при отключении от УРОВ осуществляется при включении программного ключа \textbf{«ЗАПВ от УРОВ 1С»} (\textbf{«ЗАПВ от УРОВ 2С»}).

\item
Запрет АПВ при неуспешном опробовании осуществляется при включении программного ключа \textbf{«ЗАПВ при опроб. 1С»} (\textbf{«ЗАПВ при опроб. 2С»}).

\item
Ввод оперативного запрета АПВ при отключении осуществляется внешним сигналом \textit{«Опер. запрет АПВ»}.
Оперативный запрет применяется в следующих случаях:
\begin{itemize}
\item при выполнении оперативных переключений с участием разъединителей, когда из-за возможных ошибок может сработать ДЗШ и выполнение АПВ недопустимо;
\item после проведения ремонтных работ на СШ, когда существует вероятность включения на КЗ.
\end{itemize}

\medskip

\begin{figure}[H]
\centering
\includegraphics[width=0.9\textwidth,height=0.9\textheight,keepaspectratio]{\fbpath/04.01 ДЗШ/661.9901 Логика отключения/img/Log_ZAPV.pdf}
\captionof{figure}{Функциональный блок логики запрета АПВ}\label{pic:Log_ZAPV}
\end{figure}

\medskip

\end{enumerate}


%%%%%%%%%%%%%%%%%%%%%%%%%%%%%%%%%%%%%%%%%%%%%%%%%% Логика отключения %%%%%%%%%%%%%%%%%%%%%%%%%%%%%%%%%%%%%%%%%%%%%%%%%%
\item Логика отключения

\begin{enumerate}

\item
Логика отключения предназначена для формирования сигналов отключения 1 и 2 СШ, а также сигналов отключения и запрета АПВ присоединений с учетом текущей фиксации присоединений.
Ввод и вывод функции осуществляется либо посредством программного ключа \textbf{«ЛО СШ»}, либо с использованием внешнего сигнала \textit{«Вывод ЛО СШ»}.

\item
Отключение СШ может осуществляться от ДЗШ, ЧТО, от внутренних УРОВ присоединений, а также от внешних УРОВ и внешних защит.

\item
Для фиксации срабатывания защиты на время отключения всех выключателей, время выполнения цикла АПВ, а также включения всех присоединений участвующих в этом цикле применяется сигнал \textit{«Фикс. откл. 1С»} (\textit{«Фикс. откл. 2С»}), который продлевается на выдержку времени \textbf{«Тфикс. откл. 1С»} (\textbf{«Тфикс. откл. 2С»}).
Указанный сигнал используется в логике очувствления и запрета АПВ.

\item
Логика внешнего отключения предназначена для обеспечения отключения от внешних защит (например, ЗНР ШСВ(СВ)) и для реализации действия внешних УРОВ на отключение и запрет АПВ соответствующей СШ с учетом текущей фиксации присоединения.

\item
Возможность действия внешних УРОВ на запрет АПВ задается с помощью программного ключа \textbf{«ЗАПВ внеш. Вn»}.

\item
При нарушении фиксации присоединений сигналы внешнего отключения и запрета АПВ формируются для обеих СШ.

\medskip

\begin{figure}[H]
\centering
\includegraphics[width=0.9\textwidth,height=0.9\textheight,keepaspectratio]{\fbpath/04.01 ДЗШ/661.9901 Логика отключения/img/Log_otkl_ssh.pdf}
\captionof{figure}{Функциональный блок логики отключения СШ}\label{pic:Log_otkl_ssh}
\end{figure}


\end{enumerate}

\end{enumerate}